\subsection{The Cross Product}
\begin{definition}[Cross Product]
The \textbf{cross product} is a way to multiply two 3D vectors and get a
\textbf{new vector} that is perpendicular to both of them.

\[
\vec{a} \times \vec{b}
\]

Informally, this means: \emph{take \(\vec{a}\) and \(\vec{b}\), and produce
a new vector that sticks out of the plane they form.}

Its length is

\[
|\vec{a} \times \vec{b}| = |\vec{a}| |\vec{b}| \sin \theta,
\]

which is the area of the parallelogram made by the two vectors.
\end{definition}

\subsubsection{The Right Hand Rule}
For two vectors $\vec{a}$ and $\vec{b}$, the direction of $\vec{a} \times \vec{b}$ can be determined by the right-hand rule. 

\begin{center}
\begin{tikzpicture}[scale=1.2, >=Stealth]

% Vectors
\draw[->, thick, blue] (0,0) -- (3,0) node[right] {$\vec a$};
\draw[->, thick, red] (0,0) -- (1.4,1.8) node[above] {$\vec b$};
\draw[->, thick, purple] (0,0) -- (0,2.8) node[above] {$\vec a \times \vec b$};

% Curl arrow from a to b
\draw[->, thick] (1.3,0.15) arc[start angle=8,end angle=48,radius=1.4];

% Labels
\node at (1.7,0.55) {\small curl fingers};
\node[purple] at (-0.45,1.6) {\small thumb};

% Right hand simplified
\draw[thick] (-0.25,0.15) .. controls (-0.6,0.6) and (-0.55,1.2) .. (-0.15,1.65);
\draw[thick] (-0.15,1.65) -- (0.25,1.95);
\draw[thick] (-0.10,1.35) -- (0.35,1.65);
\draw[thick] (-0.05,1.10) -- (0.45,1.35);

\end{tikzpicture}
\end{center}

\subsubsection{Magnitude of the Cross Product}
The magnitude for the cross product will depend on the 'original' vectors. 

Given vector $\vec{a}$ and $\vec{b}$, $\theta$ refers to the angle between these two vectors. 
$$\left|\vec{a} \times \vec{b}\right| = \left|\vec{a}\right|\left|\vec{b}\right|\sin \theta$$

The resultant vector can be written as:
$$\vec{r} = \left|\vec{a}\right| \left|\vec{b}\right| \sin \theta \left(\frac{1}{\left|\vec{n}\right|}\vec{n}\right)$$

Due to the cross product's non-commutative property:
$$\vec{b} \times  \vec{a} = \left|\vec{b}\right| \left|\vec{a}\right| \sin \theta (-\frac{1}{\left|\vec{n}\right|}\vec{n})$$

The magnitude of a cross-product, $\vec{u} \times \vec{v}$, turns out to be identical in value to the area of a parallelogram constructed for these vectors. 

\[
|\vec{u}\times\vec{v}|
=
|\vec{u}|\,|\vec{v}|\,\sin\theta
\]

\begin{center}
\begin{tikzpicture}[scale=1.1, >=Stealth]

\coordinate (O) at (0,0);
\coordinate (U) at (4,0);
\coordinate (V) at (1.5,2.4);
\coordinate (P) at (5.5,2.4);

\fill[blue!10] (O) -- (U) -- (P) -- (V) -- cycle;
\draw[thick] (O) -- (U) -- (P) -- (V) -- cycle;

\draw[->, very thick, blue] (O) -- (U)
    node[midway, below] {$\vec{u}$};

\draw[->, very thick, red] (O) -- (V)
    node[midway, left] {$\vec{v}$};

\draw[dashed] (V) -- (1.5,0);
\draw (1.5,0) rectangle (1.75,0.25);

\node[right] at (1.5,1.2) {$h$};

\draw[thick, ->] (0.8,0)
    arc[start angle=0,end angle=58,radius=0.8];

\node at (1.0,0.38) {$\theta$};

\end{tikzpicture}
\end{center}

\subsubsection{The Cross Product in the Cartesian Form}
\begin{definition}
    [The Cross Product]
    Given: $\vec{a} = \left[a_1, a_2, a_3\right]$ and $\vec{b} = \left[b_1, b_2, b_3\right]$
    \[
        \vec{a} \times \vec{b} = \left[a_2b_3 - a_3b_2, a_3b_1 - a_1b_3, a_1b_2 - a_2b_1\right]
    \]
\end{definition}

\begin{remark}
    When give a dot and cross product together, do the cross product first. 
\end{remark}

\subsubsection{Application of the Cross Product}
To start off, we can use the cross product to measure a torque. 
\begin{definition}
    [Moment of a force]
    \textbf{The moment of force} is the tendency of a force to tist or rotate an object about a pivot point , fulcrum or axis.
    It is a way of quantifying how much the force, which acts on the object, causes the object to rotate. 

    $$\vec{M} = \vec{r} \times \vec{f} \text{ and } \left|\vec{M}\right| = \left|\vec{r}\right| \left|\vec{f}\right| \sin \theta$$
\end{definition}

Torque is the moment of a force about a point or axis. The torque vector, $\vec{\tau}$ is defined by 
$$\vec{\tau} = \vec{r} \times \vec{f}$$